\documentclass{article}
\usepackage[utf8]{inputenc}
\usepackage{amsmath,amssymb}
\usepackage[most]{tcolorbox}
\usepackage{geometry}
\geometry{margin=2.5cm}

\newcommand{\step}[1]{\par\noindent\hangindent=3.5em\hangafter=1%
  \makebox[3.5em][l]{\textbf{#1.}}}
\newcommand{\steptag}[1]{\hfill{\small\textsf{[#1]}}}

\newtcolorbox{proofbox}[1]{
  colback=gray!5, colframe=gray!60,
  fonttitle=\bfseries, title={#1},
  breakable, enhanced,
  left=4pt, right=4pt, top=4pt, bottom=4pt,
  before skip=10pt, after skip=10pt
}

\begin{document}

\begin{proofbox}{Typed inversion derivative with chart retention: there ex...}
\step{1} Typed inversion derivative with chart retention: there exist universal C\_der, C\_ch, C\_pol, C\_grp >= 1 and kappa\_der, kappa\_ch, kappa\_pol in (0, 1/2] such that for every finite-dimensional exact-unit epsilon\_r-C*-algebra, s in \{+1, -1\}, and 0 < r <= delta satisfying C\_ch*(epsilon\_r + delta) <= kappa\_ch, C\_pol*(epsilon\_r + delta) <= kappa\_pol, C\_grp*epsilon\_r < delta - C\_pol*(epsilon\_r*delta + delta\textasciicircum{}2), C\_der*(epsilon\_r + r) <= kappa\_der, and (1 + epsilon\_r)*(1 + C\_ch*(epsilon\_r + delta))*r + C\_grp*epsilon\_r < 2*delta, the globally defined sigma(U) = u\_delta(U\textasciicircum{}dagger) maps chi\_s(B\_r\textasciicircum{}\{icalH\}(0)) into the same sJ-graph chart, where chi\_s(A) = sJ bold-dot (J + A + g\_\{sJ\}(A)), and F\_s(A) = phi\_\{sJ\}\textasciicircum{}par(sigma(chi\_s(A))) satisfies ||D(F\_s - id)(A) + 2*I\_\{icalH\}|| <= C\_der*(epsilon\_r + r) for all A in B\_r\textasciicircum{}\{icalH\}(0). \steptag{validated} \\[6pt]
\step{1.1} Constant alignment: take C\_ch,kappa\_ch from lem-stage1-unitary-graph-control; take C\_pol to be the maximum of the polar coefficients in lem-stage1-polar-retraction and lem-stage1-approximate-group-laws and kappa\_pol the minimum of their margins; and take C\_grp from lem-stage1-approximate-group-laws. Enlarging a coefficient and decreasing a margin only strengthens every hypothesis needed to invoke those lemmas. It remains to choose universal C\_der large and kappa\_der in (0,1/2] small after the universal constants in the estimates below are fixed. \steptag{validated} \\[4pt]
\step{1.2} Exact scalar-chart calculation: put V=sJ, G=g\_V, and Z\_T=T+G(T). Exact unitality, bilinearity, the involution laws in def-epsilon-cstar-algebra, and s\textasciicircum{}2=1 give f\_V(Z\_T)=G(T)+(1/2)Z\_T\textasciicircum{}dagger bold-dot Z\_T. Hence G(0)=0 by uniqueness in lem-stage1-unitary-graph-control, and, with d=C\_ch*(epsilon\_r+delta)<=kappa\_ch<=1/2, its derivative bound gives ||G(T)||<=d||T||. Differentiating G(T)=-(1/2)Z\_T\textasciicircum{}dagger bold-dot Z\_T shows ||DG(T)||<=K\_g*r whenever ||T||<=2r and epsilon\_r+r is below a universal threshold, for a universal K\_g: indeed ||DG(T)||<=(1+epsilon\_r)||Z\_T||(1+||DG(T)||), while ||Z\_T||<=(1+d)||T||<=3r, and the last DG term is absorbed. \steptag{validated} \\[4pt]
\step{1.3} Global polar factorization: for A in B\_r\textasciicircum{}\{icalH\}(0), U=chi\_s(A) belongs to calU by lem-stage1-unitary-graph-control. The common polar inverse supplied by lem-stage1-polar-retraction, identified across admissible polar data by lem-stage1-polar-coherence-naturality, and the all-calU adjoint domain in lem-stage1-approximate-group-laws give C1 maps W=sigma(U)=u\_delta(U\textasciicircum{}dagger) and Q=h\_delta(U\textasciicircum{}dagger)-J in calH with ||Q||<delta and the exact identity U\textasciicircum{}dagger=W bold-dot (J+Q). Moreover ||W-U\textasciicircum{}dagger||<=C\_grp*epsilon\_r; scalar naturality, or directly u\_delta(sJ)=sJ, also gives sigma(sJ)=sJ. \steptag{validated} \\[4pt]
\step{1.4} Polar-normal estimate: from U\textasciicircum{}dagger-W=W bold-dot Q and the preceding closeness, ||W bold-dot Q||<=C\_grp*epsilon\_r. Since W is unitary, W\textasciicircum{}dagger bold-dot W=J, and approximate associativity gives ||Q||<=||W\textasciicircum{}dagger bold-dot(W bold-dot Q)||+epsilon\_r||W||\textasciicircum{}2||Q||<=(1+epsilon\_r)||W||C\_grp*epsilon\_r+epsilon\_r||W||\textasciicircum{}2||Q||. The Cstar lower bound applied to W\textasciicircum{}dagger bold-dot W=J gives ||W||<=(1-epsilon\_r)\textasciicircum{}(-1/2). Thus, after the C\_der guard makes epsilon\_r universally small, absorption yields ||Q||<=K\_q*C\_grp*epsilon\_r for a universal K\_q. \steptag{validated} \\[4pt]
\step{1.5} Chart retention and coordinate legitimacy: the involution law and product estimate give ||U\textasciicircum{}dagger-sJ||<= (1+epsilon\_r)(1+d)||A||. Therefore ||W-sJ||<=(1+epsilon\_r)(1+C\_ch*(epsilon\_r+delta))*r+C\_grp*epsilon\_r<2delta by the stated guard. Since L\_\{sJ\}=sI exactly, C=phi\_\{sJ\}(W) has norm below 2delta; its Hermitian and anti-Hermitian projections are contractive because dagger is isometric. Writing B=C\textasciicircum{}parallel, unitarity of W implies f\_\{sJ\}(C)=0, so uniqueness in lem-stage1-unitary-graph-control forces C\textasciicircum{}perp=G(B). Consequently W=chi\_s(B) in the same sJ chart and B=F\_s(A); all maps involved are C1. \steptag{validated} \\[4pt]
\step{1.6} Exact coordinate factorization and pointwise radius: substituting U=sJ bold-dot(J+A+G(A)), W=sJ bold-dot(J+B+G(B)), and U\textasciicircum{}dagger=W bold-dot(J+Q) and cancelling the scalar s gives J-A+G(A)=(J+B+G(B)) bold-dot(J+Q). Taking the anti-Hermitian projection yields the exact identity -A=B+P\textasciicircum{}parallel((B+G(B)) bold-dot Q). Since ||G(B)||<=d||B|| and ||Q||<=K\_q*C\_grp*epsilon\_r, choosing the C\_der smallness guard so that (1+epsilon\_r)(1+d)||Q||<=1/2 gives ||B||<=r+(1/2)||B||, hence ||B||<=2r. \steptag{validated} \\[4pt]
\step{1.7} Differentiated coordinate identity: for xi in icalH set P=DF\_s(A)xi and R=DQ(A)xi. Differentiating the exact full factorization gives -xi+DG(A)xi=P+DG(B)P+R+E, where E=(P+DG(B)P) bold-dot Q+(B+G(B)) bold-dot R. Hermitian/anti-Hermitian typing gives P+xi=-P\textasciicircum{}parallel(E) and R=DG(A)xi-DG(B)P-P\textasciicircum{}perp(E). By the local graph estimate in the second step at ||A||<=r and ||B||<=2r, ||DG(A)||,||DG(B)||<=K\_g*r. \steptag{validated} \\[4pt]
\step{1.8} Uniform linear absorption: for ||xi||=1 write p=||P|| and q=||R||. The product bound, ||Q||<=K\_q*C\_grp*epsilon\_r, ||B||<=2r, ||G(B)||<=d||B||, and the local DG bounds imply ||E||<=alpha*p+beta*q with alpha<=K\_1*epsilon\_r and beta<=K\_2*r for universal K\_1,K\_2. Hence p<=1+alpha*p+beta*q and q<=K\_g*r*(1+p)+alpha*p+beta*q. If alpha,beta,K\_g*r<=1/8, elementary absorption gives p<=2, q<=K\_3*(epsilon\_r+r), and therefore ||P+xi||<=||E||<=K\_4*(epsilon\_r+r), with universal K\_3,K\_4. \steptag{validated} \\[4pt]
\step{1.9} Final constants and conclusion: choose kappa\_der>0 universally small and at most 1/2, and choose C\_der>=1 large enough, depending only on C\_grp and the universal K constants above, so that C\_der*(epsilon\_r+r)<=kappa\_der enforces every smallness and absorption condition and C\_der>=K\_4. The preceding estimate then holds uniformly for every unit xi and every A in B\_r\textasciicircum{}\{icalH\}(0): ||(DF\_s(A)+I)xi||<=C\_der*(epsilon\_r+r)||xi||. Since D(F\_s-id)+2I=DF\_s+I, this is exactly the asserted derivative inequality, while the fifth step proves the asserted same-chart retention; together with the aligned constants this proves the existential contract. \steptag{validated} \\[4pt]
\end{proofbox}

\end{document}
